\documentclass[a4paper, serif, 10pt]{moderncv}

%% moderncv
% style and color
\moderncvstyle{banking}
\moderncvcolor{black}

%% page layout/margins
\usepackage[top=2.5cm, bottom=2.5cm, left=1.5cm, right=1.5cm]{geometry}

\nopagenumbers{}

%% Babel config for Norwegian language
% ref:
% - https://latex3.github.io/babel/guides/locale-norwegian.html
\usepackage[norsk]{babel}

%% fontspec
\usepackage{fontspec}

\IfFontExistsTF{Linux Libertine}{
  \setmainfont[Ligatures={TeX, Common}, Numbers=OldStyle]{Linux Libertine}
}{}
\IfFontExistsTF{Linux Libertine O}{
  \setmainfont[Ligatures={TeX, Common}, Numbers=OldStyle]{Linux Libertine O}
}{}

%% CTeX
% CJK support, used to show CJK characters in the resume
%
% - fontset=none: disable builtin fontset but instead set the CJK font manually
% - heading=false: disable ctex heading
% - punct=kaiming: use kaiming punctuations styles for CJK
% - scheme=plain: use plain scheme, do not override \normalsize font size
% - space=auto: space settings for CJK characters
%
% ref:
% - http://ctan.mirrorcatalogs.com/language/chinese/ctex/ctex.pdf
\usepackage[UTF8, heading=false, punct=kaiming, scheme=plain, space=auto]{ctex}

\IfFontExistsTF{Noto Serif CJK SC}{
  \setCJKmainfont{Noto Serif CJK SC}
}{}
\IfFontExistsTF{Noto Sans CJK SC}{
  \setCJKsansfont{Noto Sans CJK SC}
}{}

%% line spacing
\usepackage{setspace}
\setstretch{1.125}

%% URL styling - use normal text instead of monospace
\urlstyle{same}

% Auto-underline all links
% Defer until \begin{document} so hyperref is already loaded
\AtBeginDocument{%
  \let\oldhref\href
  \renewcommand{\href}[2]{\oldhref{#1}{\underline{#2}}}%
}

%% Patch \httplink and \httpslink to handle full URLs
% The original moderncv macros always prepend http:// or https:// to URLs.
% This causes issues when the URL already has a protocol (e.g., https://example.com)
% resulting in malformed URLs like https://https://example.com.
% This patch checks if the URL already contains "://" and uses it directly if so.
% Note: We use \str_set:Nx (x-expansion) to expand macros like \@homepage first.
\ExplSyntaxOn
\str_new:N \g_yamlresume_url_str

\RenewDocumentCommand{\httpslink}{O{}m}{
  \str_set:Nx \g_yamlresume_url_str {#2}
  \str_if_in:NnTF \g_yamlresume_url_str {://}
    { \url{#2} }
    { \url{https://#2} }
}

\RenewDocumentCommand{\httplink}{O{}m}{
  \str_set:Nx \g_yamlresume_url_str {#2}
  \str_if_in:NnTF \g_yamlresume_url_str {://}
    { \url{#2} }
    { \url{http://#2} }
}
\ExplSyntaxOff

\name{Ingrid Solstad}{}
\title{Produktleder med strategisk og kommersiell erfaring}
\phone[mobile]{+47 912 34 567}
\email{ingrid.solstad@produktleder.no}
\homepage{https://www.ingridsolstad.no}

\address{Munchs gate 7, Oslo, Oslo, Norge, 0165}{}{}

\extrainfo{{\small \faLinkedin}\ \href{https://www.linkedin.com/in/ingridsolstad}{@ingridsolstad} {} {} {} • {} {} {}
{\small \faGithub}\ \href{https://github.com/ingridsolstad}{@ingridsolstad}}

\begin{document}

\maketitle

\section{Grunnleggende informasjon}

\cvline{}{Produktleder med 8 års erfaring fra digitale tjenester i finans- og teknologibransjen. Jeg jobber tett på brukere, data og forretning for å utvikle produkter som skaper reell verdi.
%
\begin{itemize}
\item Sterk på produktstrategi, veikart og prioritering
\item God til å lede tverrfaglige team og involvere interessenter
\item Erfaring med hele produktlivssyklusen – fra idé til lansering og videreutvikling
\end{itemize}}
\section{Utdanning}

\cventry{aug. 2014–juni 2016}
        {Master, Informatikk – design, bruk og interaksjon, Poeng: A (5.0)}
        {Universitetet i Oslo}
        {\url{https://www.uio.no}}
        {}
        {\begin{itemize}
\item Spesialisering i hvordan teknologi kan utformes for menneskelige behov
\item Skrev masteroppgave om brukeropplevelse i digitale selvbetjeningsløsninger
\item Gjennomførte flere prosjekter i samarbeid med norske virksomheter
\end{itemize}
\textbf{Kurs}: Interaksjonsdesign, Brukersentrert design, Prosjektledelse, Informasjonsarkitektur, Kvalitative forskningsmetoder, Datasystemer og menneskelige aktiviteter}
\section{Arbeidserfaring}

\cventry{mars 2021–Nå}
        {Senior produktleder}
        {Nordvik Bank}
        {\url{https://www.nordvikbank.no}}
        {}
        {\begin{itemize}
\item Ansvarlig for mobilbank og digitale sparingstjenester med 1,2 millioner brukere
\item Ledet produktstrategi og veikart i tett samarbeid med design, utvikling og forretning
\item Etablerte løpende brukerinnsikt og datadrevne prioriteringsprosesser
\item Økte NPS fra 48 til 66 over to år
\item Mentorerte og videreutviklet to produktmedarbeidere
\end{itemize}
\textbf{Nøkkelord}: Produktstrategi, Mobilbank, Brukerinnsikt, NPS}

\cventry{jan. 2018–feb. 2021}
        {Produktleder}
        {Aktiv Digital AS}
        {\url{https://www.aktivdigital.no}}
        {}
        {\begin{itemize}
\item Produktleder for en selvbetjeningsportal for pensjon og spareøkonomi
\item Gjennomførte brukerintervjuer, konsepttester og A/B-tester
\item Reduserte henvendelser til kundesenter med 22 prosent
\item Innførte Scrum og forbedret teamets leveransehastighet
\end{itemize}
\textbf{Nøkkelord}: Pensjon, Selvbetjening, A/B-testing, Scrum}

\cventry{aug. 2016–des. 2017}
        {Produktanalytiker}
        {Kompetansepartner AS}
        {\url{https://www.kompetansepartner.no}}
        {}
        {\begin{itemize}
\item Støttet produktleder med markedsanalyser, kundeinnsikt og konkurransescreening
\item Bidro til å definere MVP-prosjekter og forretningskrav
\item Bygget dashbord med metrikker for aktive brukere og konvertering
\end{itemize}
\textbf{Nøkkelord}: Markedsanalyse, MVP, Metrikker, Power BI}
\section{Språk}

\cvline{Norsk}{Morsmål eller tospråklig nivå \hfill \textbf{Nøkkelord}: Morsmål}
\cvline{Engelsk}{Fullt profesjonelt nivå \hfill \textbf{Nøkkelord}: Forretningsflytende}
\cvline{Svensk}{Begrenset arbeidskunnskap \hfill \textbf{Nøkkelord}: Nordisk samarbeid}
\section{Ferdigheter}

\cvline{Produktledelse}{Ekspert \hfill \textbf{Nøkkelord}: Produktstrategi, Veikart, Prioritering, OKR}
\cvline{Agile og Scrum}{Ekspert \hfill \textbf{Nøkkelord}: Scrum, Kanban, Lean Startup, Retrospektiver}
\cvline{Brukerinnsikt}{Avansert \hfill \textbf{Nøkkelord}: Intervjuer, Brukertesting, Personas, Innsiktskartlegging}
\cvline{Dataanalyse}{Viderekommen \hfill \textbf{Nøkkelord}: SQL, Google Analytics, Amplitude, A/B-testing}
\section{Utmerkelser}

\cventry{mars 2023}
        {Digital Ledelse Norge}
        {Årets produktleder 2023}
        {}
        {}
        {Hederspris for fremragende produktledelse og utviklingen av en mobilbank med tydelig brukerfokus.}
\section{Sertifikater}

\cventry{nov. 2017}
        {Scrum Alliance}
        {Certified Scrum Product Owner (CSPO)}
        {\url{https://www.scrumalliance.org/}}
        {}
        {}

\cventry{feb. 2022}
        {Google}
        {Google Analytics Individual Qualification}
        {\url{https://skillshop.exceedlms.com/}}
        {}
        {}
\section{Publikasjoner}

\cventry{sep. 2022}
        {Fra innsikt til handling – slik bygger du produkter folk vil bruke}
        {Produktledelse.no}
        {\url{https://www.produktledelse.no/artikler/fra-innsikt-til-handling}}
        {}
        {Praktisk artikkel om hvordan kvalitative og kvantitative innsikter kombineres for å forbedre produktbeslutninger.}
\section{Prosjekter}

\cventry{mars 2021–nov. 2022}
        {Ny mobilbank med økt selvbetjening og personlig økonomioversikt.}
        {Nordvik Mobil}
        {\url{https://www.nordvikbank.no/mobil}}
        {}
        {\begin{itemize}
\item Ledet produktutviklingen fra konsept til lansering
\item Etablerte tett samarbeid mellom design, utvikling og forretning
\item Økte brukertilfredsheten (NPS) med 18 poeng
\item Leverte appen i tråd med personvern- og sikkerhetskrav
\end{itemize}
\textbf{Nøkkelord}: Mobilbank, NPS, Produktlansering, Fintech}

\cventry{jan. 2018–des. 2019}
        {Selvbetjeningsportal for pensjonsinnsikt og sparing.}
        {Pensjonsportalen}
        {\url{https://www.aktivdigital.no/pensjon}}
        {}
        {\begin{itemize}
\item Gjennomførte brukerintervjuer og konsepttester
\item Utviklet enklere språk og tydeligere informasjonsarkitektur
\item Reduserte henvendelser til kundesenter med 22 prosent
\end{itemize}
\textbf{Nøkkelord}: Selvbetjening, Pensjon, Brukerinnsikt, UX}
\section{Interesser}

\cvline{Friluftsliv}{Fjellturer, Ski, Kajakk}
\cvline{Litteratur}{Faglitteratur, Teknologinytt, Podkaster}
\section{Frivillig arbeid}

\cventry{sep. 2022–Nå}
        {Karrierementor}
        {Tekna Student Oslo}
        {\url{https://www.tekna.no}}
        {}
        {\begin{itemize}
\item Veileder studenter innen teknologi og produktledelse
\item Holder workshops om veikart, prioritering og porteføljeintervjuer
\end{itemize}}

\end{document}